\section{Factoring with the same program on two amplitude representations}
Elegance is the design objective for this experiment. The ordinary Shor program uses the same modular-power and Fourier routines as a short-discrete-logarithm factoring variant. Input consists of the public modulus, the chosen base and register sizes derived from the modulus. No prime factor or order appears in circuit construction. One native modular operation multiplies by a unit modulo $N$ on labels below $N$ and fixes the remaining labels. A fixed control selects this permutation or the identity. On a GPU, inverting the multiplier with Euclid's algorithm lets each output thread read its unique input. This uses one state-vector pass per multiplication and avoids constructing a dense matrix. A fault-tolerant device must pay for reversible modular arithmetic. The costs of the two implementations have different units.

For ordinary order finding the state before the inverse Fourier transform is proportional to $\sum_a\ket a\ket{g^a\bmod N}$. For the short-log algorithm choose an odd $N=pq$ with distinct, equal-bit-length primes, with $b$ bits per prime and $m=b-1$. Define $d=(p+q-2)/2-2^m$. This is nonnegative and below $2^m$. For a suitable $g\in\mathbb{Z}_N^*$ the known element $x=g^{(N-1)/2-2^m}\bmod N$ satisfies $x=g^d$, since the difference of the exponents is $\varphi(N)/2$, a multiple of both $p-1$ and $q-1$. The order of $g$ must distinguish the desired short logarithm. Two coherent exponents produce $g^a x^{-c}$ in the work register. Fourier transformation of both exponent registers yields a frequency pair. Classical lattice post-processing recovers a candidate $d$. The two roots of $z^2-2(d+2^m+1)z+N$ give factors, each verified by exact integer arithmetic \cite{eh,shortpost,shortsuccess}.

\inputminted[firstline=100,fontsize=\footnotesize]{qaf}{examples/prefix/factoring.qaf}
The register roles in this listing are explicit. The first exponent has $m+\ell$ wires and the second has $\ell$ wires. The work register begins at their combined width. Base-two order finding modulo 21 used fifteen wires and recovered order six, giving $3\cdot7$. The balanced short-log instance modulo 143 used seventeen wires with $m=\ell=3$. It recovered $d=3$, yielding $11\cdot13$. The first sampled pair was $(51,5)$. All sixteen seeded samples in each small experiment yielded verified factors with the selected post-processing parameters. These frequencies come from the simulated circuit. The independent reference enumerates the public modular function on its exponent registers and Fourier transforms each work-label indicator. It agrees with every marginal probability to within $6.2\times10^{-16}$ on the two examples. The reference supplies no period to the execution. Quaspy provides the classical post-processing routines \cite{orderpost,quaspy}.

As of 8 November 2026, optimizing RSA factoring depends on a physical model and a cost objective. Gidney combines approximate residue arithmetic with fault-tolerant resource choices \cite{gidney}. The Pinnacle architecture considers quantum LDPC codes, and reconfigurable-atom analyses explore different time--space tradeoffs \cite{pinnacle,oratomic}. Pappalardo and colleagues' October cat-qubit preprint gives one estimate of 22,593 physical qubits and 3.14 days for RSA-2048 under a one-microsecond error-correction cycle, ten-microsecond reaction, specified phase-flip noise and negligible bit flips \cite{catshor}. A different operating point in that work spends more qubits for a shorter schedule. These are projections. Our demonstrations implement the logical Fourier and modular computations on a classical GPU. They do not implement the resource papers' physical codes, magic-state factories, layouts or RSA-2048 schedules. The general classical simulation still scales exponentially.

\subsection{When the whole network denotes a qubit}
Let $F_\theta:\{0,1\}\to\C$ be the output of an entire network and require $0<\sum_b|F_\theta(b)|^2<\infty$. Its denotation is the ray of $F_\theta/\|F_\theta\|$. An embedding with two complex basis weights represents every such ray in the ideal parameter model. Define $\theta\sim\theta'$ when these denotations coincide. Mapping an equivalence class to its ray is injective by definition and surjective onto the reachable rays. \texttt{Networks.lean} proves this fact for any preparation map. For this architecture its image is all pure qubits. For raw weights the preparation is many-to-one: $(1,0)$ and $(2,0)$ give the same qubit. A register with $n$ inputs initially represents a product of normalized two-label functions. As a circuit runs, the backend evaluates the resulting joint function on $\B^n$.

We executed the same seventeen-wire prefix file under native CUDA and a PyTorch amplitude backend. In PyTorch the input wire was an instance of \texttt{BasisAmplitudes} initialized with trainable weights $(2,0)$, whose normalized vector is $\ket0$. The existing first $X$ and every subsequent modular and Fourier operation were taken directly from the source. Comparing all $131{,}072$ output amplitudes gave maximum absolute difference $5.21\times10^{-18}$ or less. The check compares complex amplitudes, including relative phases, before sampling. It establishes equivalence on this run. On entangled intermediate states the representation is a joint amplitude computation.

The same language can declare a general feedforward network directly. This network supplies the real and imaginary parts of a spatial amplitude.
\inputminted{qaf}{examples/prefix/wave_net.qaf}
The exporter produces ordinary runnable PyTorch \cite{torch}. It uses the same declared layer ordering as the live \texttt{nn.Sequential} model. The example loads identical parameters in the generated and live modules and compares their outputs.
\inputminted[fontsize=\footnotesize]{python}{examples/wave_net.py}
For the box problem set $\psi_j\propto x_j(1-x_j)(u_\theta(x_j)+i v_\theta(x_j))$ on 40 interior nodes. The optimizer differentiates the Rayleigh quotient of the actual tridiagonal Dirichlet Hamiltonian on the GPU. It reduced the energy from $5.48168220$ to $4.93238870$. The independent grid eigenvalue is $4.93238821$, while the continuum ground energy is $4.93480220$. The first discrepancy measures variational error at the selected grid size, and the second comparison includes spatial discretization. A uniform approximation theorem for continuous functions would not by itself certify either error. Energy convergence for an unbounded Hamiltonian requires control in its form domain.

\subsection{One build, two representations}
\texttt{build\_qaf.py} takes one circuit file and creates an executable CUDA plan together with a Lean \texttt{Program.lean}. The build also contains generated CUDA expressions, the device kernels, the input state and gate tables, and a manifest with content hashes. The verify step reloads the exported plan, compiles its saved CUDA kernels and replays it without parsing QAF. It also evaluates the Lean program and checks that its recovered UTF-8 bytes equal the original. For the factoring program CUDA replay was bit-identical to direct native execution. The Lean source isomorphism is a theorem over the embedding. CUDA replay is an execution check. A separate Lean theorem identifies the portable instruction datatype with its supported Quil-line datatype in both directions. A full proof about floating-point execution, GPU drivers or a provider optimizer remains open.
